Met het \texttt{sudo}\index{sudo} (switch user do\index{switch user do}) commando kan je commando's uitvoeren alsof ze van root, of een andere gebruiker, zijn, je moet daar dan natuurlijk wel de rechten voor hebben niet iedereen mag dat doen. Wie die rechten heeft wordt bepaald door het \texttt{/etc/sudoers} bestand of door bestanden in de \texttt{/etc/sudoers.d} directory of, zoals we eerder gezien hebben, door lid te zijn van de sudo-groep.

Om een commando uit te voeren als root, hoeven we net als bij \texttt{su} aan \texttt{sudo} niets mee te geven. We gaan dit testen door de home-directory van root te bekijken:
\begin{lstlisting}[language=bash]
$ sudo ls /root
\end{lstlisting}
Hiermee hebben we aangetoond dat we commando's als root kunnen uitvoeren.

Aan \texttt{sudo} kunnen we ook de \texttt{-u} optie meegeven om een andere gebruiker te zijn:
\begin{lstlisting}[language=bash]
$ sudo -u mies ls /home/mies
\end{lstlisting}
Natuurlijk hebben we geen \texttt{mies} gebruiker op ons systeem, dus dit commando zal een error geven.

